\section{The arithmetic orientation and finite torsors}
\label{sec:do26-orientation}

The finite arithmetic Chern--Simons invariant uses two maps.  A torsor
pulls a group-cohomology class to an arithmetic curve.  An arithmetic trace
then evaluates that class.  The first map is a construction in the category
of sheaves.  The second uses class field theory.  We separate these maps
and keep the primes dividing the coefficient order in the base scheme.

Let \(F\) be a totally imaginary number field, fix \(n\geq1\),
and put \(X=\operatorname{Spec}\mathcal O_F\).  Throughout this section,
\(H^q(X,-)\) means ordinary sheaf cohomology on the small \'etale site
\(X_{\mathrm{et}}\).  Its objects are schemes \'etale over \(X\), and
its coverings are jointly surjective \'etale families.  Coefficients are
sheaves of abelian groups, with cohomological grading.  In particular,
\(\mathbb G_m(U)=\Gamma(U,\mathcal O_U)^\times\), and
\(\underline{\mathbb Z/n}\) is the constant sheaf.  No completion or
inverse limit of coefficient sheaves occurs.  We impose no compact-support
condition and use no compactified site in this construction.

\subsection{Roots of unity and the power-map defect}

For a field \(L\) of characteristic zero, put
\(\mu_n(L)=\{u\in L^\times:u^n=1\}\).
Suppose \(\mu_n(\overline F)\subset F\), and choose
\(\zeta\in F\) of exact order \(n\).  For \(n=1\), this means
\(\zeta=1\).  The corresponding group isomorphism is
\[
 \zeta_n:\tfrac1n\mathbb Z/\mathbb Z\longrightarrow\mu_n(F),
 \qquad k/n\longmapsto\zeta^k.
\]
Write \(\mu_{n,X}=\ker([n]:\mathbb G_m\to\mathbb G_m)\), where
\([n](u)=u^n\).  This is a sheaf on the specified small site.

\begin{lemma}[The coefficient map at every finite prime]
\label{lem:do26-roots}
The map
\[
 e_\zeta:\underline{\mathbb Z/n}\longrightarrow\mu_{n,X},
 \qquad k\longmapsto\zeta^k,
\]
is an isomorphism.  Moreover,
\[
 \operatorname{Hom}_{X_{\mathrm{et}}}
   (\underline{\mathbb Z/n},\mathbb G_m)=\mu_n(F).
\]
\end{lemma}

\begin{proof}
A root of unity is an algebraic integer, and its inverse is an algebraic
integer.  Thus \(\zeta\in\mathcal O_F^\times\), so the displayed map
is defined on all of \(X\).  At a geometric point over a finite prime,
the stalk of \(\mathbb G_m\) is
\((\mathcal O_{X,\bar x}^{\mathrm{sh}})^\times\).
The strict henselization is a discrete valuation domain of characteristic
zero.  Its fraction field contains \(F\), and
\[
 T^n-1=\prod_{k=0}^{n-1}(T-\zeta^k)
\]
has exactly those \(n\) distinct roots in that field.  The same argument
applies to the geometric generic stalk.  Therefore \(e_\zeta\) is
bijective on every geometric stalk, which proves the first assertion.
A homomorphism from the constant sheaf \(\underline{\mathbb Z}\) to
any abelian sheaf is determined by the image of its global section \(1\).
Factoring through \(\underline{\mathbb Z/n}\) requires that image to be
killed by \(n\).  For \(\mathbb G_m\), these sections are precisely
\(\mu_n(\mathcal O_F)=\mu_n(F)\).
\end{proof}

The stalk in this proof is a strict henselian local ring.  It is not the
geometric special fiber of the group scheme \(\mu_n\).  At a prime
dividing \(n\), that group scheme need not be \'etale.

For example, take \(F=\mathbb Q(i)\), \(n=2\), and the prime
\((1-i)\).  In the strict henselian local ring, normalize the valuation
by \(v(1-i)=1\).  The unit \(i\) has no square root there.  Indeed,
\(u^2=i\) would imply \(\bar u=1\) in the residue field of
characteristic two.  Both \(u-1\) and \(u+1\) would have positive
valuation, whereas
\[
  v\bigl((u-1)(u+1)\bigr)=v(i-1)=1.
\]
This is impossible.  Thus the \'etale power map can fail to be surjective,
even under the root-of-unity hypothesis.

Define its image and cokernel in the category of abelian sheaves:
\[
 K_n=\operatorname{im}([n]:\mathbb G_m\to\mathbb G_m),
 \qquad Q_n=\operatorname{coker}([n]:\mathbb G_m\to\mathbb G_m).
\]
The cokernel is a sheaf cokernel.  Its sections need not equal the quotient
of the groups of sections.

\begin{lemma}[Nonzero defect at every coefficient prime]
\label{lem:coefficient-prime-defect}
Let \(v\) be a finite prime of residue characteristic \(p\mid n\).
The stalk of \(Q_n\) at a geometric point over \(v\) is nonzero.
This assertion holds for every number field, without a root-of-unity
hypothesis.
\end{lemma}

\begin{proof}
Let \(R\) be the strict henselization at that point and choose a
uniformizer \(\varpi\), with valuation \(v(\varpi)=1\).
The unit \(1+\varpi\) is not a \(p\)-th power in \(R\).
Indeed, if \(x^p=1+\varpi\), reduction to the residue field gives
\(\bar x=1\), so \(x=1+a\) with \(v(a)\geq1\). In
\[
 (1+a)^p-1=\sum_{k=1}^{p-1}\binom pk a^k+a^p,
\]
every summand has valuation at least two: \(p\) divides each
intermediate binomial coefficient and \(p\geq2\).
This contradicts \(v(\varpi)=1\). An \(n\)-th power is a
\(p\)-th power, so \(1+\varpi\) represents a nonzero class in
\(R^\times/(R^\times)^n\), the stalk of \(Q_n\).
\end{proof}

\begin{lemma}[Unramified norms on units]
\label{lem:do26-norm}
Let \(R\) be a henselian discrete valuation ring with finite residue
field \(k\).  Let \(R'/R\) be its unramified extension of degree
\(d\geq1\).  Then
\(N_{R'/R}:R'^\times\to R^\times\) is surjective.
\end{lemma}

\begin{proof}
Write \(|k|=q\).  The norm on residue fields sends
\(x\in\mathbb F_{q^d}^\times\) to
\(x^{(q^d-1)/(q-1)}\).  Since that multiplicative group is cyclic,
the image has order \(q-1\).  It is therefore all of \(k^\times\).

Choose an \(R\)-basis of the free module \(R'\).  In those coordinates,
the norm is the determinant of multiplication and hence a polynomial.
Fix \(a\in R^\times\), and choose a residue unit \(\bar b\) with
norm \(\bar a\).  The derivative of the norm at \(\bar b\) is
\[
 h\longmapsto N(\bar b)
       \operatorname{Tr}_{\mathbb F_{q^d}/\mathbb F_q}(h/\bar b).
\]
The trace is a nonzero linear functional because the finite-field
extension is separable.  Some coordinate derivative is therefore nonzero.
Lift all coordinates of \(\bar b\) to \(R\), and hold all but that
coordinate fixed.  Hensel's lemma solves the resulting one-variable norm
equation with the prescribed residue.  The solution is a unit of \(R'\)
with norm \(a\).
\end{proof}

\begin{lemma}[Cohomology of the power-map defect]
\label{lem:do26-defect}
For every number field \(F\) and every \(n\geq1\),
\[
 \operatorname{supp}(Q_n)=\{v:v\mid n\},
 \qquad H^q(X,Q_n)=0\quad(q\geq1).
\]
These conclusions do not require roots of unity in \(F\).
\end{lemma}

\begin{proof}
If \(n\) is invertible at a geometric point, the equation
\(T^n-u=0\) for a unit \(u\) has a simple root after a finite
separable residue extension.  Hensel's lemma lifts that root.  Thus the
stalk of \(Q_n\) vanishes away from the finite set \(V(n)\).
Lemma~\ref{lem:coefficient-prime-defect} proves nonvanishing at
each point of that set; for \(n=1\), both sides are empty.
For its reduced closed immersion \(i:V(n)\hookrightarrow X\), the
adjunction map \(Q_n\to i_*i^*Q_n\) is an isomorphism on every stalk.
Closed direct image is exact on the \'etale site and preserves injectives,
since its inverse-image left adjoint is exact.  Consequently it suffices
to compute cohomology over each finite residue field.

Fix a prime \(v\mid n\).  Put \(R=\mathcal O_{X,v}^{\mathrm h}\)
and let \(R^{\mathrm{sh}}\) be a strict henselization with residue
field \(\bar k\).  The relevant discrete
\(G_k=\operatorname{Gal}(\bar k/k)\)-module is
\[
 M=(R^{\mathrm{sh}})^\times/
                \bigl((R^{\mathrm{sh}})^\times\bigr)^n.
\]
It is killed by \(n\).  Write \(\phi\) for arithmetic Frobenius,
which generates the dense cyclic subgroup of \(G_k\cong\widehat{\mathbb Z}\).
For a discrete torsion module over this group,
\[
 H^1(G_k,M)=M/(\phi-1)M,\qquad H^q(G_k,M)=0\quad(q\geq2).
\]
Here cochains are continuous.  The first formula follows by evaluating
a cocycle on \(\phi\).  An element with orbit length \(r\) and order
dividing \(e\) gives a cocycle on the cyclic quotient of order \(re\).
For the second formula, use the two-periodic cyclic resolution and take
the limit over finite quotients.  If the action factors through a cyclic
quotient of order \(r\), inflation to the quotient of order \(rs\)
has a factor \(s^{\lfloor q/2\rfloor}\) in degree \(q\).
Choosing \(s\) divisible by the exponent kills each class with
\(q\geq2\).  Every continuous cochain has finite image and lies in
such a finite torsion submodule.

We prove \((\phi-1)M=M\).  Represent \(m\in M\) by a unit
\(u\in R_r^\times\), where \(R_r/R\) is finite unramified of
degree \(r\).  Choose the unramified extension \(R_{rn}/R_r\)
of degree \(n\), and put \(\sigma=\phi^r\) on it.  By
Lemma~\ref{lem:do26-norm}, there is a unit \(b\) with
\(N(b)=u\).  Therefore
\[
 z=u b^{-n}\quad\text{satisfies}\quad N(z)=1.
\]
The elementary cyclic form of Hilbert's theorem 90 gives
\(z=\sigma(w)/w\) in the fraction field of \(R_{rn}\).
For completeness, choose \(t\) for which
\[
 w=\sum_{j=0}^{n-1}
    \left(\prod_{\ell=0}^{j-1}\sigma^\ell(z^{-1})\right)\sigma^j(t)
 \ne0.
\]
Such a \(t\) exists by the linear independence of distinct field
automorphisms.  The equation \(N(z)=1\) makes
\(\sigma(w)=zw\) by cyclic reindexing.  Since the extension is
unramified, a uniformizer of \(R_r\) has valuation one upstairs.
Multiplying \(w\) by a suitable power of that uniformizer makes it a
unit and does not change \(\sigma(w)/w\).
In the additive notation of \(M\), we obtain
\[
 [u]=(\phi^r-1)[w]
     =(\phi-1)(1+\phi+\cdots+\phi^{r-1})[w].
\]
This proves the required surjectivity and all higher-cohomology vanishings.
\end{proof}

\subsection{The global trace and its normalization}

\begin{lemma}[The unramified local invariant]
\label{lem:do26-local-invariant}
Let \(R\) be a henselian discrete valuation ring with finite residue
field \(k\), fraction field \(L\), and normalized valuation \(v\).
Let \(L^{\mathrm{nr}}=\operatorname{Frac}(R^{\mathrm{sh}})\) and
\(G_k=\operatorname{Gal}(L^{\mathrm{nr}}/L)\).
Valuation, the inverse Bockstein, and evaluation at arithmetic Frobenius
\(\phi\) give an isomorphism
\[
 H^2(G_k,(L^{\mathrm{nr}})^\times)
 \xrightarrow{v_*}H^2(G_k,\mathbb Z)
 \xrightarrow{\beta^{-1}}H^1(G_k,\mathbb Q/\mathbb Z)
 \xrightarrow{\operatorname{ev}_{\phi}}\mathbb Q/\mathbb Z.
\]
Here \(\beta\) is the boundary for
\(0\to\mathbb Z\to\mathbb Q\to\mathbb Q/\mathbb Z\to0\).
\end{lemma}

\begin{proof}
Put \(U=(R^{\mathrm{sh}})^\times\).  On a finite unramified cyclic
extension, the odd positive cohomology of its units is zero by Hilbert's
theorem 90.  Its proof for units follows from the explicit formula and
valuation adjustment in Lemma~\ref{lem:do26-defect}.  The even positive
cohomology is zero by the cyclic resolution and
Lemma~\ref{lem:do26-norm}.  Taking the limit over finite unramified
extensions gives \(H^q(G_k,U)=0\) for every \(q>0\).
The valuation sequence
\(0\to U\to(L^{\mathrm{nr}})^\times\xrightarrow{v}\mathbb Z\to0\)
therefore makes \(v_*\) an isomorphism in degree two.
Positive-degree cohomology of the trivial module \(\mathbb Q\)
vanishes by averaging over finite quotients.  Thus \(\beta\) is an
isomorphism.  Finally, every continuous character of
\(\widehat{\mathbb Z}\) is determined by its value on \(\phi\),
and every element of \(\mathbb Q/\mathbb Z\) occurs.
\end{proof}

The preceding local calculation isolates the global arithmetic input.
We use the following class-field trace theorem on the ordinary small site:
\begin{equation}
\label{eq:do26-trace-input}
 H^2(X,\mathbb G_m)=0,\qquad
 \operatorname{tr}_F:H^3(X,\mathbb G_m)
                   \xrightarrow{\sim}\mathbb Q/\mathbb Z.
\end{equation}
The field is totally imaginary in this statement.  These are the
class-field computations in Mazur~\cite[\S2,
(2.1)--(2.3), pp.~537--538]{Mazur1973}.  We specify how their normalization enters.

Write \(\operatorname{Br}(L)=H^2(L,\mathbb G_m)\) for the
cohomological Brauer group of a field \(L\).
Let \(j:\operatorname{Spec}F\to X\) be the generic-point morphism and
let \(i_v:\operatorname{Spec}k_v\to X\) run through the finite primes.
The divisor sequence is
\[
 0\longrightarrow\mathbb G_m\longrightarrow j_*\mathbb G_{m,F}
   \xrightarrow{(v)_v}\bigoplus_v i_{v*}\underline{\mathbb Z}
   \longrightarrow0,
\]
where \(v\) takes a local uniformizer to \(+1\).
Local class field theory gives \(R^qj_*\mathbb G_{m,F}=0\) for
\(q>0\).  The resulting cohomology sequence contains
\begin{equation}
\label{eq:do26-residue}
\begin{split}
0\longrightarrow H^2(X,\mathbb G_m)&\longrightarrow\operatorname{Br}(F)
 \xrightarrow{\partial}\bigoplus_v\mathbb Q/\mathbb Z\\
 &\xrightarrow{\delta}H^3(X,\mathbb G_m)
 \longrightarrow H^3(F,\mathbb G_m)\longrightarrow0.
\end{split}
\end{equation}
The copy of \(\mathbb Q/\mathbb Z\) at \(v\) is
\(H^2(k_v,\underline{\mathbb Z})\), identified using arithmetic
Frobenius and the boundary of
\(0\to\mathbb Z\to\mathbb Q\to\mathbb Q/\mathbb Z\to0\).
The invariant on unramified Brauer classes is
the map of Lemma~\ref{lem:do26-local-invariant}.  Its sign can be
computed without a convention for central simple algebras.  For an
unramified cyclic extension of degree \(d\), with arithmetic Frobenius
generator \(\sigma\) and uniformizer \(\varpi\), the cocycle
\[
 \alpha(\sigma^i,\sigma^j)=
          \varpi^{\lfloor(i+j)/d\rfloor},\qquad 0\leq i,j<d,
\]
has invariant \(1/d\).  Indeed, its valuation is the Bockstein of the
character \(\sigma\mapsto1/d\), using the lift
\(\sigma^i\mapsto i/d\) to \(\mathbb Q\).

The substantive class-field assertions are that \(\partial\) agrees
with these local invariants, that
\[
 0\longrightarrow\operatorname{Br}(F)
  \xrightarrow{\partial}\bigoplus_v\mathbb Q/\mathbb Z
  \xrightarrow{\sum_v}\mathbb Q/\mathbb Z\longrightarrow0
\]
is exact, and that \(H^3(F,\mathbb G_m)=0\).
They imply~\eqref{eq:do26-trace-input}, with the normalization
\begin{equation}
\label{eq:do26-trace-normalization}
 \operatorname{tr}_F\bigl(\delta((a_v)_v)\bigr)=\sum_v a_v.
\end{equation}
The connecting homomorphism in~\eqref{eq:do26-residue} uses the usual
cochain differential: lift a cocycle, apply the differential, and take its
class in the kernel.  Formula~\eqref{eq:do26-trace-normalization} fixes
the sign.  The local unit argument does not prove these global class-field
assertions.  They remain the arithmetic hypotheses whenever the trace
theorem is not invoked.

\begin{proposition}[The degree-three orientation]
\label{prop:do26-orientation}
Assume the trace theorem~\eqref{eq:do26-trace-input}, with
normalization~\eqref{eq:do26-trace-normalization}.  Inclusion of roots of
unity induces an isomorphism
\[
 H^3(X,\mu_{n,X})\xrightarrow{\sim}
       H^3(X,\mathbb G_m)[n].
\]
Consequently
\begin{equation}
\label{eq:do26-orientation}
 \operatorname{inv}_{F,n,\zeta}
 =\operatorname{tr}_F\circ H^3(\mu_{n,X}\hookrightarrow\mathbb G_m)
                 \circ H^3(e_\zeta)
 :H^3(X,\underline{\mathbb Z/n})
       \xrightarrow{\sim}\tfrac1n\mathbb Z/\mathbb Z.
\end{equation}
If \(u\in(\mathbb Z/n)^\times\), then
\(\operatorname{inv}_{F,n,\zeta^u}
   =u\operatorname{inv}_{F,n,\zeta}\).
\end{proposition}

\begin{proof}
The definition of \(K_n\) gives two short exact sequences:
\[
 0\to\mu_{n,X}\to\mathbb G_m\xrightarrow{[n]}K_n\to0,
 \qquad
 0\to K_n\xrightarrow{\iota}\mathbb G_m\to Q_n\to0.
\]
Lemma~\ref{lem:do26-defect} and the second sequence show that
\(H^r(\iota)\) is an isomorphism for \(r\geq2\).  In particular,
\(H^2(X,K_n)=0\).  The first sequence therefore identifies
\(H^3(X,\mu_{n,X})\) with the kernel of
\(H^3(X,\mathbb G_m)\to H^3(X,K_n)\).
After composing with \(H^3(\iota)\), this map is multiplication by
\(n\).  The trace identifies its kernel with
\((\mathbb Q/\mathbb Z)[n]=\tfrac1n\mathbb Z/\mathbb Z\).
Lemma~\ref{lem:do26-roots} supplies the remaining isomorphism.
Finally, \(e_{\zeta^u}=e_\zeta\circ[u]\).  Functoriality of
cohomology gives the stated multiplier, with no inversion of \(u\).
\end{proof}

\begin{corollary}[The top-degree duality pairing]
\label{cor:do26-duality}
Under the same trace hypothesis, the Yoneda evaluation pairing
\[
 H^3(X,\underline{\mathbb Z/n})\times
 \operatorname{Hom}_X(\underline{\mathbb Z/n},\mathbb G_m)
 \longrightarrow\mathbb Q/\mathbb Z,
 \qquad(x,f)\longmapsto\operatorname{tr}_F(f_*x),
\]
is perfect.  If \(f(1)=\zeta^a\), its value is
\(a\operatorname{inv}_{F,n,\zeta}(x)\).
\end{corollary}

\begin{proof}
The formula follows from \(f=(\mu_{n,X}\hookrightarrow\mathbb G_m)
\circ e_\zeta\circ[a]\).  Proposition~\ref{prop:do26-orientation}
identifies the first factor with \(\tfrac1n\mathbb Z/\mathbb Z\).
Lemma~\ref{lem:do26-roots} identifies the second with the cyclic group
generated by \(\zeta\).  The resulting pairing \((t,a)\mapsto at\)
identifies each finite group with the full character group of the other.
For \(n=1\), both groups are zero and this remains valid.
\end{proof}

This proves the required degree-three part from the trace theorem.  It
does not assume perfectness of the full Artin--Verdier pairing in every
degree.  Such perfectness is the stronger theorem in
Mazur~\cite[\S2, (2.4), pp.~538--539]{Mazur1973}.

\subsection{The classifying pullback}

Let \(A\) be a finite discrete group.  A right \(A\)-torsor on
\(X_{\mathrm{et}}\) is a sheaf \(P\) with a right action, locally
isomorphic to \(\underline A\) acting on itself.  For a finite constant
group it is represented by a finite \'etale scheme over \(X\).
Fix a geometric generic point \(b:\operatorname{Spec}\overline F\to X\)
and put \(\pi=\pi_1^{\mathrm{et}}(X,b)\).
The finite-cover equivalence identifies finite \'etale schemes with finite
continuous left \(\pi\)-sets
\cite[Theorem 58.6.2]{do26-StacksFundamental}.

Given a continuous homomorphism \(\rho:\pi\to A\), give the set
\(A\) the action \(\gamma\cdot a=\rho(\gamma)a\).
This action commutes with right multiplication.  The finite-cover
equivalence therefore produces a right torsor \(P_\rho\).
Conversely, a point \(p\in P_b\) determines \(\rho\) by
\(\gamma p=p\rho(\gamma)\).  Replacing \(p\) by \(pa\) replaces
\(\rho\) by \(a^{-1}\rho a\).  The equivalence preserves morphisms,
so it also identifies torsor isomorphisms with these changes of point.
Thus
\[
 \{\text{finite \'etale right }A\text{-torsors on }X\}/\cong
 =\operatorname{Hom}_{\mathrm{cont}}(\pi,A)/A.
\]
No surjectivity of \(\rho\) is required.  A torsor need not be connected.

Let \(BA\) denote the classifying topos of left \(A\)-sets.
For a left \(A\)-set \(M\), define the associated sheaf
\[
 T_P(M)=P\times^A\underline M,
 \qquad (pa,m)\sim(p,am),
\]
where the quotient is formed as a sheaf quotient.  Locally on a cover
trivializing \(P\), this is the constant-sheaf functor on the underlying
set of \(M\).  It preserves finite limits and all colimits.  Its right
adjoint sends a sheaf \(\mathcal F\) to
\(\operatorname{Hom}_X(P,\mathcal F)\), with action
\((a\cdot f)(p)=f(pa)\).  Hence \(T_P\) is the inverse image of a
geometric morphism \(f_P:X_{\mathrm{et}}\to BA\).

\begin{proposition}[Cohomology pulled back by a torsor]
\label{prop:do26-pullback}
Let \(M\) be a discrete abelian group with trivial \(A\)-action.
For each \(q\geq0\), a right \(A\)-torsor \(P\) determines a
homomorphism
\[
 f_P^*:H^q(A,M)\longrightarrow H^q(X,\underline M).
\]
It depends only on the isomorphism class of \(P\).
\end{proposition}

\begin{proof}
On abelian group objects the associated-sheaf functor
\[
 T_P:\mathbb Z[A]\text{-}\mathrm{Mod}
       \longrightarrow\mathrm{Ab}(X_{\mathrm{et}})
\]
is exact.  Exactness can be checked after trivializing \(P\), where
it is the exact constant-sheaf functor.  Since the action on \(M\) is
trivial, \([p,m]\mapsto m\) identifies \(T_P(M)\) canonically with
\(\underline M\).  Also \(T_P(\mathbb Z)=\underline{\mathbb Z}\).

Use the Yoneda descriptions
\[
 H^q(A,M)=\operatorname{Ext}^q_{\mathbb Z[A]}(\mathbb Z,M),
 \qquad
 H^q(X,\underline M)=
 \operatorname{Ext}^q_{X}(\underline{\mathbb Z},\underline M).
\]
Apply \(T_P\) to an exact sequence representing a Yoneda extension.
Exactness gives an extension on \(X\) with the stated endpoints.
The functor preserves equivalences of extensions and their Baer sums.
It therefore gives the displayed homomorphism.  In degree zero it sends
an invariant element to its constant section.
A torsor isomorphism \(P\to P'\) induces the natural isomorphism
\([p,m]\mapsto[p',m]\) of these exact functors.  On the trivial
endpoints this is the identity, so the resulting extension classes agree.
\end{proof}

For \(P=P_\rho\), write this homomorphism as \(\rho^*\).
The notation includes the passage from group cohomology to sheaf
cohomology.  It asserts no isomorphism
\(H^q(\pi,M)\cong H^q(X,\underline M)\), and requires no
\(K(\pi,1)\) property for \(X\).

\subsection{Conjugation, cocycles, and the invariant}

We make the conjugation identity explicit.  For trivial coefficients,
inhomogeneous cochains have differential
\[
\begin{split}
 (dc)(g_1,\ldots,g_{q+1})={}&c(g_2,\ldots,g_{q+1})\\
 &+\sum_{i=1}^{q}(-1)^i
       c(g_1,\ldots,g_i g_{i+1},\ldots,g_{q+1})\\
 &+(-1)^{q+1}c(g_1,\ldots,g_q).
\end{split}
\]
For a profinite group and a finite coefficient group, cochains are
continuous maps to the discrete coefficient group.  For finite \(A\),
all maps are continuous.

\begin{lemma}[An explicit conjugation homotopy]
\label{lem:do26-conjugation}
Let \(c\in Z^3(A,M)\), with trivial \(A\)-action on \(M\), and
let \(a\in A\).  Set \(\operatorname{Ad}_a(g)=aga^{-1}\) and
\[
\begin{split}
 h_a(c)(g,h)={}&c(a^{-1},\operatorname{Ad}_a g,
                         \operatorname{Ad}_a h)\\
              &-c(g,a^{-1},\operatorname{Ad}_a h)
                +c(g,h,a^{-1}).
\end{split}
\]
Then \(d h_a(c)=c\circ\operatorname{Ad}_a-c\).
\end{lemma}

\begin{proof}
Use homogeneous bar chains, with differential
\[
 \partial(g_0,\ldots,g_q)=
 \sum_{j=0}^{q}(-1)^j(g_0,\ldots,\widehat g_j,\ldots,g_q).
\]
Here \(q\geq1\), and \(\partial_0=0\).
The chains carry the diagonal left \(A\)-action.  Right translation
\(R(g_0,\ldots,g_q)=(g_0a^{-1},\ldots,g_qa^{-1})\) commutes
with this action.  Define
\[
 K_q(g_0,\ldots,g_q)=
 \sum_{i=0}^{q}(-1)^i
 (g_0,\ldots,g_i,g_i a^{-1},\ldots,g_q a^{-1}).
\]
Deleting a vertex away from the repeated index cancels against the
corresponding term of \(K_{q-1}\partial\).  At the repeated index,
successive terms cancel.  The two endpoints give
\[
 \partial K_q+K_{q-1}\partial=R-\operatorname{id}.
\]
The convention in degree zero has \(K_{-1}=0\).

Let \(\widetilde c\) be the homogeneous cochain associated to \(c\):
\[
 \widetilde c(g_0,\ldots,g_q)=
 c(g_0^{-1}g_1,\ldots,g_{q-1}^{-1}g_q).
\]
It is left invariant.  This invariance identifies \(R^*\widetilde c\)
with the homogeneous cochain of \(c\circ\operatorname{Ad}_a\).
Dualizing the chain identity gives
\(d(\widetilde c K_2)=R^*\widetilde c-\widetilde c\), since
\(d\widetilde c=0\).  Evaluating \(K_2\) at \((1,g,gh)\)
gives exactly the three displayed terms of \(h_a(c)\).
\end{proof}

\begin{theorem}[The finite arithmetic Chern--Simons function]
\label{thm:do26-kim}\label{thm:kim}
Let \(F\), \(X\), \(n\), and \(\zeta\) satisfy the hypotheses of
Lemma~\ref{lem:do26-roots}.  Use the class-field trace
\eqref{eq:do26-trace-input} with normalization
\eqref{eq:do26-trace-normalization}.  For a finite group \(A\) and a
class \(c\in H^3(A,\mathbb Z/n)\) with trivial coefficient action,
\[
 \operatorname{CS}_{c,\zeta}:
 \operatorname{Hom}_{\mathrm{cont}}(\pi_1^{\mathrm{et}}(X,b),A)/A
   \longrightarrow\tfrac1n\mathbb Z/\mathbb Z,
 \qquad [\rho]\longmapsto
       \operatorname{inv}_{F,n,\zeta}(\rho^*c),
\]
is well defined. For \(u\in(\mathbb Z/n)^\times\), replacing
\(\zeta\) by \(\zeta^u\) multiplies the function by \(u\).
\end{theorem}

\begin{proof}
Proposition~\ref{prop:do26-pullback} constructs \(\rho^*c\) in the
required cohomology group.  It is independent of the cocycle representing
\(c\), because the construction is defined on extension classes.
If \(\rho'=a\rho a^{-1}\), the map of fibers \(x\mapsto ax\)
is \(\pi\)-equivariant and commutes with the right \(A\)-action.
It therefore induces an isomorphism \(P_\rho\to P_{\rho'}\).
The two pullbacks agree by Proposition~\ref{prop:do26-pullback}.
At the group-cochain level, Lemma~\ref{lem:do26-conjugation} gives
the same invariance by an explicit coboundary.
Proposition~\ref{prop:do26-orientation} gives the value and its
dependence on the root.  A change of geometric base point gives an
equivalence of finite-cover fiber functors, unique up to inner
conjugation, so it does not alter the function on torsor isomorphism
classes.  This is the finite construction of
Kim~\cite[\S1, pp.~1--3]{Kim2016}.
\end{proof}

For \(A=1\), positive-degree group cohomology vanishes, so this function
is zero.  It is also zero for \(n=1\).  Neither conclusion requires
\(H^3(X,\underline{\mathbb Z/n})\) to vanish for other values of \(n\).

\subsection{The ramification boundary}

The group \(\pi_1^{\mathrm{et}}(X,b)\) classifies extensions
unramified at every finite prime, including primes above \(n\).
For a finite set \(S\) of finite primes, put
\(X_S=\operatorname{Spec}\mathcal O_F[S^{-1}]\).
A representation of \(\pi_1^{\mathrm{et}}(X_S,b)\) gives the
pullback of Proposition~\ref{prop:do26-pullback} on \(X_S\).
An extension to a torsor on \(X\) requires trivial inertia at every
prime of \(S\).  To see this criterion, normalize \(X\) in the
finite generic algebra of the torsor.  Over a discrete valuation ring,
trivial inertia is equivalent to an unramified extension.  Its residue
extensions are separable because finite fields are perfect.  The
normalization is then finite \'etale at those primes and at all primes
of \(X_S\), and hence supplies the extension.

Kim's boundary construction takes \(S\) to contain every prime above
\(n\). For \(v\in S\), let \(F_v\) be the completion of \(F\)
and let \(G_{F_v}=\operatorname{Gal}(\overline F_v/F_v)\). Choose
embeddings that give restriction maps
\(G_{F_v}\to\pi_1^{\mathrm{et}}(X_S,b)\). Given a local
representation \(\rho_v\) and a cocycle \(c\), the set
\[
 \mathcal H_v(\rho_v)=
 \{\beta\in C^2_{\mathrm{cont}}(G_{F_v},\mathbb Z/n):
                      d\beta=\rho_v^*c\}/B^2
\]
where \(B^2=dC^1_{\mathrm{cont}}(G_{F_v},\mathbb Z/n)\),
is either empty or a torsor under
\(H^2(G_{F_v},\mathbb Z/n)\).  Indeed, the difference of two
solutions is a two-cocycle, and quotienting by coboundaries gives a
free transitive action.  Nonemptiness follows from the local vanishing
\(H^3(F_v,\mathbb Z/n)=0\), proved with the invariant
identification in Lemma~\ref{lem:do27-local}. The chosen root
identifies the acting group with \(\tfrac1n\mathbb Z/\mathbb Z\).
A finite product of these torsors can then be pushed out by the
sum of local invariants, as in
Kim~\cite[\S2, pp.~4--6]{Kim2016}.
Those local primitives and their compatibility are additional data.
The construction on \(X\) above provides no scalar evaluation of an
arbitrarily ramified representation on \(X_S\).

\subsection{Infinite places}

The totally imaginary hypothesis excludes real decomposition groups.
To discuss other fields, let \(F\) now be an arbitrary number field,
put \(X=\operatorname{Spec}\mathcal O_F\), and let \(X_\infty\)
be the set of real embeddings and conjugate pairs of complex embeddings.
The notation \(\overline{\operatorname{Spec}\mathcal O_F}\) alone
does not select a cohomology theory.

The Artin--Verdier site \(\widetilde X_{\mathrm{et}}\) has objects
\((Y,M)\), where \(Y\to X\) is separated and \'etale,
\(M\subset Y_\infty\), and a complex point of \(M\) does not map
to a real point of \(X_\infty\). A morphism carries the chosen
infinite subset into the chosen infinite subset. A family covers when
it jointly covers both the schemes and those subsets. Coefficients on
this site mean specified sheaves, together with their infinite-place
specialization data; they are not determined by a topological
compactification \cite[\S1]{Bienenfeld1987}.

At a real place the decomposition group is \(C_2\). For the trivial
module \(\mathbb Z/2\), its two-periodic resolution has zero
cochain differentials, so
\[
 H^q(C_2,\mathbb Z/2)=\mathbb Z/2\qquad(q\geq0).
\]
For odd-order coefficients, the higher groups vanish by averaging.
Thus ordinary real-place cohomology has no finite
\(2\)-cohomological dimension. A three-dimensional duality statement
with real places must specify its framework: ordinary small-site
cohomology with the \(2\)-primary part excluded; modified cohomology
with Tate cohomology at the real places; or derived sheaf cohomology on
an explicitly defined extended Artin--Verdier site, with a constructible
torsion sheaf, its extension, and its infinite-place specialization
maps. The latter two are different constructions. The trace used in
Proposition~\ref{prop:do26-orientation} has the stated totally imaginary
ordinary small-site domain \cite[\S2]{Mazur1973}.
